\addchap{A Anmerkung zur Nutzung von Künstlicher Intelligenz}
\setcounter{chapter}{1}

Im Rahmen dieser Arbeit wurden Künstliche Intelligenz (\acrshort{acr:KI})\index{Künstliche Intelligenz} basierte Werkzeuge benutzt. Tabelle~\ref{tab:anhang_uebersicht_KI_werkzeuge} gibt eine Übersicht über die verwendeten Werkzeuge und den jeweiligen Einsatzzweck.

\begin{table}[hbt]	
	\centering
	\renewcommand{\arraystretch}{1.5}
	\captionabove[Liste der verwendeten KI-basierten Werkzeuge]{Liste der verwendeten \acrshort{acr:KI}-basierten Werkzeuge}
	\label{tab:anhang_uebersicht_KI_werkzeuge}
	\begin{tabular}{>{\raggedright\arraybackslash}p{0.25\linewidth} >{\raggedright\arraybackslash}p{0.70\linewidth}}
		\textbf{Werkzeug} & \textbf{Beschreibung der Nutzung}\\
		\hline 
		\hline
		Claude (Anthropic) & \vspace{-\topsep}
					\begin{itemize}[noitemsep,topsep=0pt,partopsep=0pt,parsep=0pt] 
						\item Unterstützung bei der Code-Modularisierung und Refactoring
						\item Überprüfung der LaTeX-Dokumentation
						\item Recherche zu ESP32-S3 GPIO-Einschränkungen und Sensor-Datenblättern
						\item Unterstützung bei der LVGL~9 Migration und UI-Entwicklung
					\end{itemize} \\
		\hline 
	\end{tabular} 
\end{table}

Sämtliche durch KI-Werkzeuge generierten oder überarbeiteten Inhalte wurden von den Autoren geprüft, verifiziert und bei Bedarf angepasst. Die inhaltliche Verantwortung liegt ausschließlich bei den Autoren.

%=============================================================================
\addchap{B Kostenaufstellung}
\setcounter{chapter}{2}
\setcounter{section}{0}
\setcounter{table}{0}
%=============================================================================

Die Gesamtkosten für das InspectAir-Projekt betragen ca.\ \textbf{138,88\,\euro{}}. Tabelle~\ref{tab:kosten_detail} zeigt die detaillierte Aufstellung aller beschafften Komponenten.

\textbf{Wichtiger Hinweis zu den Stückkosten:} Die meisten Kleinteile (Kondensatoren, Widerstände, Buttons, Kabel, Steckverbinder) wurden in Multipacks mit mehreren Stück pro Packung bestellt. Die tatsächlichen Stückkosten pro verwendetem Bauteil sind daher deutlich geringer als die angegebenen Paketpreise. Der 3D-Druck des PETG-Gehäuses war \textbf{kostenlos} verfügbar.

\begin{table}[htb]
\centering
\renewcommand{\arraystretch}{1.2}
\captionabove[Detaillierte Kostenaufstellung]{Detaillierte Kostenaufstellung aller Komponenten}
\label{tab:kosten_detail}
\small
\begin{tabular}{p{7cm} r r}
\toprule
\rowcolor{tableheader}
\textbf{Bauteil} & \textbf{Kosten} & \textbf{Kategorie} \\
\midrule
\multicolumn{3}{l}{\textit{Hauptsensoren}} \\
MH-Z19C CO\textsubscript{2}-Sensor & 17,05\,\euro{} & Sensor \\
Plantower PMS5003 (Feinstaub) & 13,69\,\euro{} & Sensor \\
GY-SGP40 TVOC Sensor & 6,68\,\euro{} & Sensor \\
AHT20 (Temperatur/Feuchte) & 1,19\,\euro{} & Sensor \\
Hi-Link HLK-LD2410C Radar & 3,49\,\euro{} & Sensor \\
\midrule
\multicolumn{3}{l}{\textit{Mikrocontroller und Display}} \\
Waveshare ESP32-S3-Devkit N8R8 WROOM & 17,85\,\euro{} & MCU \\
Display 4,0 Zoll TFT (ST7796S) & 15,92\,\euro{} & Display \\
\midrule
\multicolumn{3}{l}{\textit{Platinen und Verkabelung}} \\
Lochrasterplatine & 5,99\,\euro{} & Platine \\
Breadboard (groß) & 9,99\,\euro{} & Prototyping \\
Breadboard (klein) & 1,58\,\euro{} & Prototyping \\
Zubehör für Platine (Buchsenleisten, Draht, Kabel) & 10,53\,\euro{} & Zubehör \\
Jump Wires & 1,64\,\euro{} & Kabel \\
Jumperkabel-Set & 2,47\,\euro{} & Kabel \\
Adapter für PMS5003 & 9,35\,\euro{} & Adapter \\
\midrule
\multicolumn{3}{l}{\textit{Elektronische Bauteile (Multipacks)}} \\
I2C Logic Level Shifter & 1,52\,\euro{} & IC \\
Widerstandsset & 2,64\,\euro{} & Passiv \\
Keramikkondensatoren & 0,99\,\euro{} & Passiv \\
Elektrolytkondensatoren (220\,µF 25V) & 0,99\,\euro{} & Passiv \\
Sicherung PTVC & 1,12\,\euro{} & Schutz \\
\midrule
\multicolumn{3}{l}{\textit{Bedienelemente und Montage}} \\
Buttons RA & 0,99\,\euro{} & Taster \\
Buttons normal & 1,35\,\euro{} & Taster \\
USB C Panel Mount & 6,99\,\euro{} & Stecker \\
Schrauben und Muttern & 2,97\,\euro{} & Montage \\
\midrule
\multicolumn{3}{l}{\textit{3D-Druck}} \\
PETG-Gehäuse (3D-Druck) & \textbf{0,00\,\euro{}} & Gehäuse \\
\midrule
\midrule
\textbf{Gesamtsumme} & \textbf{135,99\,\euro{}} & \\
\bottomrule
\end{tabular}
\end{table}

\paragraph{Zusammenfassung nach Kategorien:}
\begin{itemize}[nosep]
\item \textbf{Sensoren:} ca.\ 42\,\euro{} (5 Sensoren)
\item \textbf{MCU + Display:} ca.\ 34\,\euro{}
\item \textbf{Platinen \& Verkabelung:} ca.\ 42\,\euro{}
\item \textbf{Elektronische Bauteile:} ca.\ 7\,\euro{} (Multipacks)
\item \textbf{Bedienelemente \& Montage:} ca.\ 12\,\euro{}
\item \textbf{3D-Druck:} kostenlos
\end{itemize}

Bei einer Serienproduktion würden die Stückkosten durch Wegfall der Multipack-Aufschläge und direkten Komponenteneinkauf auf geschätzt \textbf{80--90\,\euro{}} sinken.

%=============================================================================
\addchap{C Pin-Belegung und Bauteilliste}
\setcounter{chapter}{3}
\setcounter{section}{0}
\setcounter{table}{0}
\setcounter{figure}{0}
%=============================================================================

Die vollständige Pin-Belegung und Bauteilliste liegt als separates Dokument (\textit{Pin\_Belegung\_Bauteilliste.pdf}) vor. Nachfolgend eine Zusammenfassung der kritischen Verbindungen.

\section{Kritische Hinweise}

\textbf{LEVEL SHIFTER PFLICHT:} Der LD2410C sendet TX-Signale mit 5\,V-Pegel. Der ESP32 verträgt maximal 3,3\,V. Ohne Level Shifter wird der ESP32 zerstört. Der Level Shifter muss zwischen LD2410C TX und ESP32 GPIO\,7 geschaltet werden.

\textbf{STÜTZKONDENSATOREN PFLICHT:} MH-Z19C und LD2410C verursachen Stromspitzen. 220\,µF (25\,V) Elkos direkt an VCC/GND der Sensoren!

\textbf{KERAMIKKONDENSATOREN:} 100\,nF Keramikkondensatoren an 3,3\,V VCC des Displays und 3,3\,V VCC des ESP32 zur Entstörung.

\textbf{I2C PULL-UPS:} 4,7\,k$\Omega$ Widerstände von SDA (GPIO\,8) und SCL (GPIO\,18) nach 3,3\,V.

%=============================================================================
\addchap{D Übersicht der Engineering-Dokumente}
\setcounter{chapter}{4}
\setcounter{section}{0}
\setcounter{table}{0}
%=============================================================================

Das Engineering Folder umfasst folgende Dokumente, die als separate PDFs beiliegen:

\begin{table}[hbt]
\centering
\renewcommand{\arraystretch}{1.3}
\captionabove[Engineering Folder]{Übersicht der Engineering-Dokumente}
\label{tab:eng_folder}
\begin{tabular}{r p{4cm} p{7cm}}
\toprule
\rowcolor{tableheader}
\textbf{\#} & \textbf{Dokument} & \textbf{Inhalt} \\
\midrule
1 & Requirements Specification & Funktionale, nicht-funktionale und Hardware-Anforderungen \\
2 & Design Description & System- und Software-Architektur \\
3 & Coding Rules & Namenskonventionen, Doxygen, Git-Workflow \\
4 & Configuration Management & Entwicklungsumgebung, Libraries, Pinout \\
5 & Schedule & Zeitplan mit 7 Meilensteinen \\
6 & Risk \& Opportunities & Risikoanalyse und Erweiterungsmöglichkeiten \\
7 & Testspezifikation & 28 Testfälle mit Pass/Fail-Kriterien \\
8 & Pin-Belegung \& BOM & GPIO-Zuordnung und Bauteilliste \\
\bottomrule
\end{tabular}
\end{table}
